\section{Space of curve functionals}

\subsection{Topology}

Geodesic currents are positive Radon measures on $UT\Sigma$, and as such they can be seen, by Riesz representation theorem, as linear positive functionals on the space of compacty supported functions on $UT\Sigma$.
However, this viewpoint is forgetting the fact that currents are also measures invariant under the geodesic flow.
In this paper, we incorporate this observation and consider geodesic currents as functionals on curves $\F(S)$, i.e., functions $f \colon \Curves(S) \to \mathbb{R}$.
The space of $\Curves(S)$ is just thought, a priori, as a set, so these are just real valued functionals.
However, we can endow $\F(S)$ with a topology which doesn't require any topology on $\Curves(S)$ to define.

Indeed, given a geodesic current $\mu$, we can consider the function intersection number $i(\mu, \cdot) \colon \Curves(S) \to \mathbb{R}$. 

Work of Otal \cite[Th\`eor\'eme~2]{Otal90:SpectreMarqueNegative} shows that the map
\[
\GC(S) \hookrightarrow \F(S), \mu \mapsto \left( i(\mu,\cdot) \colon \Curves(S) \to \mathbb{R} \right).
\]
is injective.

From this, we can show, as in \cite[Theorem~10]{DLR10:DegenerationFlatMetrics}, the following Proposition.

\begin{proposition}
Let $(\mu_i)_{i \in \mathbb{N}}$ be a sequence with $\mu_i \in \GC(S)$, and $\mu \in \GC(S)$. Then,
$\mu_i \to \mu$ in the weak$^*$ topology convergence of geodesic currents if and only if $i(\mu_i,C) \to i(\mu,C)$ for all closed curve $C \in \Curves(S)$.
\end{proposition}

Following this, we define the following topology on functionals. 

\section{Pointwise convergence}

If $\F$ is a family of functionals defined on a set $X$ to a topological space $Y$, then $\F$ is contained in the product 
\[
Y^X \coloneqq \prod \{ Y \mid x \in X \}.
\]
The topology $\mathcal{P}$ of \emph{pointwise convergence} is the following.
A net $\{ f_n \colon n \in D\}$ converges to $g$ if and only if $\{ f_n(x) \colon n \in D\}$ converges to $g(x)$ for each $x \in X$. A subbase for $\mathcal{P}$ is given by the family of all subsets of the form $\{ f \colon f(x) \in U \}$, where $x \in X$ and $U$ is open in $Y$. For each point $x\in X$ there is a function \emph{evalution at } $x$, denoted $\epsilon_x \in \F$ and defined by $\epsilon_x(f)=f(x)$ for all $f \in \F$.
Evaluation at $x$ is continuous and open with respect to the topology $\mathcal{P}$ \cite[Theorem~3.2]{K75:GenTop}, and $\mathcal{P}$ is the smallest topology for $\F$ such that evaluation is continuous. A function $g$ on a topological space to $\F$ is continuous with respect to $\mathcal{P}$ if and only if $\epsilon_x \circ g$ is continuous for each $x \in X$ \cite[Theorem~3.3]{K75:GenTop}. Observe that the topology on $X$ does not play a role in the definition of $\mathcal{P}$, only the family of functionals $\F$ and the topology of $Y$ play a role. If $Y$ is Hausdorff then $\F$ is also Hausdorff \cite[3.5]{K75:GenTop}.
A family $\F$ is \emph{poitnwise closed} iff $\F$ is a closed subspace of the product space $Y^X$.
If $A$ is a subset of $X$ then $\F[A]$ is the set of all points $f(x)$ with $x \in A$.
If $x \in X$, then $\F[\{x\}]$ will be abbreviated to $\F[x]$. Note that $\epsilon_x(F)=F[x]$.

If $Y$ is Hausdorff, we have the following characterization of compactness, proven in \cite[Chapter~7,Theorem~1]{K75:GenTop}

\begin{theorem}
 $\F$ is compact in the topology $\mathcal{P}$ if and only if
 \begin{enumerate}
     \item $\F$ is pointwise closed in $Y^X$, and
     \item for each point $x \in X$, the set $F[x]$
 \end{enumerate}
 \label{thm:pointwisecompact}
\end{theorem}

We will use the topology of pointwise convergence on functionals on curves.

\begin{definition}
Let $(f_i)$ be a net with $f_i \in \F(S)$, and $f \in \F(S)$.  Then
 $f_i \to f$ if and only if $f_i(C) \to f(C)$ for all closed curve $C \in \Curves(S)$.
\end{definition}




\begin{definition}
  An internal current $\lambda \in \GC_0(\Sigma)$ is \emph{filling} if every geodesic in $G(\Sigma)$ which does not project to $\partial \Sigma$ is transversely intersected by a geodesic in the support of $\lambda$. We denote by $\GC_{\mbox{fill}}(S) \subset \GC(S)$ be the set of all filling currents.
\end{definition}

\begin{example}
\begin{enumerate}
    \item For $S$ a compact surface possibly with non-empty boundary, let $C \in C(\Sigma) \subset \GC(\Sigma)$ be a closed curve so that $S - C$ is a collection of disks and annuli, where each annuli contains a component of $\partial S$. This is an example of a filling current which is not of full support.
    \item The Lioville current is an example of filling current of full support.
\end{enumerate}
\end{example}

Given a surface $S$ of finite type, a subsurface $S_v \subset S$ and a geodesic current $\mu \in \GC(S)$, the \emph{systole of $\mu$ on $S_v$} is defined as
\[
  \sys_{S_v}(\mu) \coloneqq \inf_{c \in \Curves(S_v-\partial S_v)} i(\mu,c)
\]
The systole of $\mu$ will be denoted by $\sys(\mu)$ and defined to be $\sys(\mu) \coloneqq \sys_{S}(\mu)$.

There are several characterization of filling currents in terms of intersection number for currents.

\begin{proposition}
   Let $\mu \in \GC_0(S)$.
   The following are equivalent:
   \begin{enumerate}
       \item $\sys(\mu)>0$.
       \item $i(\mu,\eta)>0$ for all $\eta \in \GC_0(S)$ for $\eta \neq 0$.
       \item $i(\mu,\lambda) > 0$ for all $\lambda\in \ML(S)$.
   \end{enumerate}
   \label{prop:filling}
\end{proposition}
\begin{proof}
  See \cite[Theorem~4.1]{BIPP17:Structure}.
\end{proof}

\begin{proposition}
  The space of filling currents is dense and open subset in $\GC_0(S)$
\end{proposition}
\begin{proof}
  See \cite[Lemma~3.6]{EM18:ErgodicGeodesicCurrents} for denseness and \cite[Corollary~3.8]{EM18:ErgodicGeodesicCurrents} for openness.
\end{proof}

\section{Uniform boundedness principle for currents?}

We want to prove that pointwise convergence of convex function implies convergence on compact sets.

Idea: 

1) extend functional to signed currents.
2) hit it with functional analysis (uniform boundedness). Issue, to do this, we need a Banach space, and therefore need the norm topology on signed currents, which is stronger than the weak* topology.

So, it seems we can't use functional analysis to prove this.
On the other hand, in order to say anything interesting about these functionals we should probably work with norm topology.




Let $X$ be the space of geodesic currents. Let $f \colon X \to \mathbb{R}$ be a continuous functional satisfying $f(\alpha + \beta) = f(\alpha) + f(\beta)$, and $f(a \alpha) = a f(\alpha)$ for $a \geq 0$.
We will call such a functional $f$ \emph{linear}.
We define a distance in geodesic currents by
\[
d(x,y) \coloneqq \sum_{\alpha \in \Curves(S) \subset X} a_{\alpha} |i(x,\alpha)-i(y,\alpha)|
\]
where $i(\cdot,\cdot)$ denotes the intersection number on geodesic currents, and the real positive coefficients $a_{\alpha}$ are conveniently chosen so that the sum is finite.
We define the \emph{norm} of a geodesic current by
\[
\|x\| \coloneqq d(x,0).
\]
Finally, we define the norm of a linear functional $f$ by
\[
\| f \| \coloneqq \sup_{\| x \| \leq 1} \frac{|f(x)|}{\|x \|}.
\]
Given a linear functional on geodesic currents, we can extend $f$ to signed geodesic currents, by defining $\bar{f}(\mu) \coloneqq f(\mu^+) - f(\mu^-)$ where $\mu^+,\mu^-$ are the unique non-negative measures obtained from the Hahn decomposition of $\mu$. 

\begin{remark}We warn that it doesn't follow in general that $\bar{f}$ is continuous on signed currents. In fact, it is not necessarily true that if a net $\mu_i \to \mu$, then $\mu^+_i \to \mu^+$ and $\mu^-_i \to \mu^-$.
However, we have the following Proposition.
\end{remark}

\begin{proposition}
Let $\mu_i,\mu$ be signed measures, and suppose $\mu_i \to \mu$ and $|\mu_i|(X) \to |\mu|(X)$, where $|\mu|\coloneqq \mu^+ + \mu^-$.
Then $\mu^+_i \to \mu^+$ and $\mu^-_i \to \mu^-$.
\end{proposition}
\begin{proof}
See \cite[Corollary~8.4.8]{Bo07:MeasureTheory}.
\end{proof}

\begin{theorem}
Let $\F$ be a family of linear functionals, and suppose $\sup_{f \in \F} |f(x)| < \infty$ for any $x \in X$. Then $\sup_{f \in \F} \| f\| < \infty$.
\label{thm:boundedness}
\end{theorem}
\begin{proof}
For every $n \in \mathbb{N}$, then we define
\[
X_n \coloneqq \{ x \in X \colon \sup_{f \in \F} |f(x)| \leq n \}.
\]
Note thast $X \coloneqq \cup_n X_n$.
$X$ is a complete metric space. Thus, by the Baire category theorem, there exist an $m>0$ so that $X_m$ has non-empty interior, and thus there exists $x_0 \in X_m$ and $\epsilon>0$, so that
\[
B_{\epsilon}(x_0) \coloneqq \{ x\in X \colon d(x,x_0) \leq \epsilon \} \subset X_m.
\]
Now, for any $u \in X$, with $\| u \| \leq 1$ we have
\begin{align*}
|f(u)|=  \epsilon^{-1} |\bar{f}(x_0 + \epsilon u - x_0)| & \leq \epsilon^{-1}|\bar{f}(x_0 + \epsilon u) - \bar{f}(x_0)| \leq \\
\epsilon^{-1}(|f(x_0 + \epsilon u)| + |f(x_0)|) \leq 2\epsilon^{-1}m
\end{align*}
In the third equality we used the definition of the extension $\bar{f}$.
Taking supremum over all $\| u \| \leq 1$, we obtain the desired result.
\end{proof}

We say that a functional $f \colon X \times X \to \mathbb{R}$ is \emph{bilinear} if $f(\cdot,y) \colon X \to \mathbb{R}$ and $f(x,\cdot) \colon X \to \mathbb{R}$ are linear for each $x,y \in X$.
We say that a functional $f \colon X \times X \to \mathbb{R}$ is \emph{separate continuous} if $f(\cdot,y) \colon X \to \mathbb{R}$ and $f(x,\cdot) \colon X \to \mathbb{R}$ are continuous for each $x,y \in X$.

\begin{theorem}
Let $g \colon X \times X \to \mathbb{R}$ be separate continuous and bilinear.
Then $g$ is continuous.
\end{theorem}
\begin{proof}
Let $(x_n,y_n) \to (0,0)$.
Let $U$ be a neighborhood of $0$.
Since $X$ is locally compact, we can assume $U$ is closed ball, hence compact.
For each $n$, the function $f_n \colon g(x_n, \cdot )$ is continuous.
For all large enough $n$, $x_n \in U$.
Thus, for each $y \in U$, we have $\{ f_n(y) \colon n \in \mathbb{N} \}$ is bounded, since $U$ is compact and $g(\cdot,y)$ is continuous on $U$.
Thus, by Theorem~\ref{thm:boundedness}, we have
\[
  M\coloneqq \sup_{n \in \mathbb{N}} \| f_n \| < \infty.
\]
Therefore,
\[
  |g(x_n,y_n)| = |f_n(y_n)| \leq \|f_n \| \| y_n\| \leq M \|y_n\| \to 0
\]
The first equality uses the definition of $f_n$, the second inequality the definition of norm of $\| f_n \|$, the third inequality the uniform bound and the last step uses the fact $y_n \to 0$.
If $(x_n,y_n) \to (x,y)$ we use bilinearity of $f$ to note that
$y_n' \coloneqq y_n-x$ is a signed current, and moreover $|y_n'|(X) \to 2y(X)$ and $(y_n')^+ \to (y')^+$ and $(y_n')^- \to (y')^-$, therefore $\bar{f_n}$ is continuous. Thus, for any $\epsilon>0$, there exist large enough $n_0$ so that for all $n \geq n_0$,
\[
  |\bar{g}(x_n-x,y_n-y)| \leq |\bar{f_n}(y_n-y)| + |\bar{g}(x,y_n-y)| \leq \|f_n \| \| y_n-y\| + \epsilon \leq M \|y_n-y\| + \epsilon \leq (M+1)\epsilon.
\]
\end{proof}




\section{Geodesic currents as measures}


Taking $X=UT\Sigma$, we can see geodesic currents $\GC(\Sigma)$ as non-negative finite Radon measures on $X$ invariant under the geodesic flow. They sit inside the space of \emph{signed geodesic currents} $\hat{\GC}(\Sigma)$, signed finite Radon measures on $X$ invariant under the geodesic flow.

\begin{claim}
The Kantorovich-Rubinstein norm $\| \|$ on $\Meas^+(X)$ defines a norm on the space of signed geodesic currents $\hat{\GC}(\Sigma)$.
\end{claim}
\begin{proof}
The same definition of $\| \|$ applies to geodesic flow invariant measures?
\end{proof}

From this Claim and Proposition~\ref{prop:convexcont}, it follows that if a convex function $f \colon \GC(S) \to \mathbb{R}$ is locally upper bounded in some neighborhood of $\mu \in \GC(S)$, then it is automatically continuous.


\section{Bounds on continuous, positive, homogeneous functionals}

\begin{definition}[Positive]
  We say a functional $f \colon \GC(S) \to \mathbb{R}$ is positive if $f(\mu)>0$ for all $\mu \in \GC(S)$, where $\mu \neq 0$.
\end{definition}

We note the following.

\begin{proposition}
A convex functional is non-negative.
\end{proposition}

This motivates the following definition of filling convex functional.

\begin{definition}
A \emph{filling} functional is a convex functional $f$ so that its extension to geodesic currents is positive.
\end{definition}

The following result will be of use.

\begin{theorem}
Given $\mathcal{K}$ a compact subset of $\F_{\mbox{fill}}(S)$, there is a constant $K>1$ depending only on $\mathcal{K}$ so that
\[
\frac{1}{K}\cdot f_1(\mu) \leq f_2(\mu) \leq K \cdot f_1(\mu).
\]
for all $f_1,f_2 \in \mathcal{K}$ and any $\mu \in \GC(S)$.
\label{thm:lipfillfunc}
\end{theorem}
\begin{proof}
Consider the function
\[
F \colon \mathcal{K} \times \mathcal{K} \times \GC(S) \to \mathbb{R}, (f_1,f_2,\mu) \mapsto \frac{f_1(\mu)}{f_2(\mu)}.
\]
Since $f_1,f_2$ are filling, $F$ is well-defined and continuous, since evaluation at $\mu$ by $f_1$ and $f_2$ is continuous.
Moreover, since $f_1,f_2$ are homogeneous, $F$ descends to a continuous function
\[
F \colon \mathcal{K} \times \mathcal{K} \times \mathbb{P}\GC(S) \to \mathbb{R}
\]
Since $\mathcal{K} \times \mathcal{K} \times \mathbb{P}\GC(S)$ is compact, it follows that $F$ is bounded from below and above by positive numbers.
\end{proof}

Two specializations will be of use.

\begin{corollary}
Let $\Sigma$ be a compact hyperbolic possibly with non-empty boundary.
Given any two continuous, homogeneous and positive functionals $f_1,f_2 : \GC(\Sigma) \to \mathbb{R}$, then
there exist a positive constant $K$ so that
\[
   \frac{1}{K}\cdot f_2 \leq f_1 \leq K \cdot f_2.
\]
\label{prop:funclip}
\end{corollary}

\begin{corollary}
 Let $K \subset S$ be a compact subset, and $A \subset \GC_{\mbox{fill}}(S)$ be compact sets. There exists a positive constant $R \geq 1$ (depending on $K$ and $A$, so that
 \[
    \frac{1}{R} \cdot  i(\lambda_1,\mu) \leq i(\lambda_2,\mu) \leq R \cdot i(\lambda_1,\mu)
 \]
 for all $\mu \in \GC_K(S)$ and all $\lambda_1,\lambda_2 \in A$.
\end{corollary}

\section{Filling functionals and conjectures}

\begin{lemma}
A convex functional $f$ is filling if $f(\lambda) > 0$ for any measured lamination $\lambda \in \ML$.
\end{lemma}
\begin{proof}
If $f$ is filling, we clearly have the result, since $\ML(S)\subset \GC(S)$.
Suppose now that $f(\lambda)>0$ for all $\lambda \in \ML$.
Fix $\mu$. Take any sequence $(a_i C_i)_{i \in \mathbb{N}}$ with $a_i \geq 0$ real weights, and $C_i$ closed curves, and $a_i C_i \to \mu$. Let $\bar{C_i}$ denote the simple multi-curve obtained after smoothing all essential crossings from $C_i$.
By the smoothing property, we have
\[
  f(C_i) \geq f(\bar{C_i})
\]
for all $i$.
By homogeneity
\[
  f(a_iC_i) \geq f(a_i\bar{C_i})
\]
and, in the limit, we get, by continuity of $f$,
\[
  f(\mu) \geq f(\lambda)
\]
where $\lambda = \lim_i a_i\bar{C_i}$, with $\lambda \in \ML$, since $i(\lambda,\lambda)=a_i^2i(C_i,C_i)=0$.
Therefore, $f(\mu) >0$.
\end{proof}



Denote $\F_{\mbox{fill}}(S) \subset \F(S)$ the subset of filling convex functionals.
We have the following conjectures.

\begin{conjecture}
A convex functional $f$ is filling if and only if
$\sys(f)>0$.
\label{conj:filling_systole}
\end{conjecture}
\begin{proof}
Assume $f$ is filling. Showing $\sys(f)>0$ is easy.
Apply Proposition~\ref{prop:funclip} with $f_1=f$ and $f_2=i(\mathcal{L}_{\Sigma},\cdot)$, where $\Sigma$ is a fixed choice of hyperbolic metric. We get then a constant $K^f_{\Sigma}$ so that
\[
  \frac{1}{K^f_{\Sigma}}\cdot i(\mathcal{L}_{\Sigma},C) \leq f(C) \leq K^f_{\Sigma}\cdot i(\mathcal{L}_{\Sigma},C).
\]
Thus, $\sys(f)>0$.
Now, assume $\sys(f)>0$. This requires something else. Analyze proof of geodesic current case?
\end{proof}

\begin{conjecture}
A convex filling functional $f$ has finitely many systoles.
\end{conjecture}
\begin{proof}
A convex filling functional is positive, same idea as in the Proof of Conjecture~\ref{conj:filling_systole} applies. Then the result follows since 

\begin{align*}
  \{ \mbox{ connected C } \in \Curves(S) \mid f(C) = \sys(f) \} &\subset  \{ \mbox{ connected C } \in \Curves(S) \mid f(C) \leq \sys(f) \} \subset \\
  \mbox{ connected C } \in \Curves(S) \mid \frac{1}{K^f_{\Sigma}} \ell_{\Sigma}(C) \leq \sys(f) \}
\end{align*}
and the cardinality of the last set is finite.
\end{proof}

Let $\phi \in \MCG$, and $f \in \F(S)$.
We define the action of $\MCG$ as $\phi \cdot f(C) \coloneqq f(\phi \cdot C)$ for any curve $C \in \Curves(S)$.


\begin{conjecture}
$\MCG$ acts properly discontinuously on the space of filling functionals $\F_{\mbox{fill}}(S)$, i.e., the set
\[
 \{ \phi \in \MCG \mid \phi(K) \cap K \neq \emptyset \}
\]
is finite for all non-empty compact $K \subset \F_{\mbox{fill}}(S)$.
\label{conj:propdisc}
\end{conjecture}

First we have the following Lemma \cite[Lemma~4.3]{EM18:ErgodicGeodesicCurrents}.

\begin{lemma}
Let $C \in \Curves(S)$ be a multicurve which is filling. Then the set
\[
\{ \phi \in \MCG \mid i(C,\phi(C)) \leq L \}
\]
is finite for all $L \geq 0$.
\label{lem:finitefilling}
\end{lemma}

The proof of the Conjecture will follow from Lemma~\ref{lem:finitefilling} and a similar argument as  in \cite[Lemma~4.5]{EM18:ErgodicGeodesicCurrents}.

\begin{proof}[Proof of Conjecture~\label{conj:propdisc}]
Let $\mathcal{K} \subset \F_{\mbox{fill}}(S)$, and let $(\phi_n)$ be a sequence of pairwise distinct elements in $\MCG$.
It suffices to show $\mathcal{K} \cap \phi_n\mathcal{K} = \emptyset$ for all large $n$.
Fix a filling multicurve $C$, and let
\[
m_n = \min \{ f(\phi_n^{-1}(C)) \mid f \in \mathcal{K} \}
\]
for each $n$.
By Theorem~\ref{thm:pointwisecompact}, the subspace $\mathcal{K}[\phi_n^{-1}(C)]=\{ f(\phi_n^{-1}(C)) \mid f \in \mathcal{K} \}$ is compact. Thus, its minimum is achieved and so $m_n >0$, since $f \in \mathcal{K}$ are filling functionals.
It also follows that $\max \{ f(\phi_n^{-1}(C)) \mid f \in \mathcal{K} \} < \infty$.
Therefore, it is enough to show that $m_n \to \infty$ as $n \to \infty$.
By  Proposition~\ref{thm:lipfillfunc}, with $f_1=f$ and $f_2=i(C,\cdot)$, using that $C$ is filling, we get a constant $K \geq 1$ so that
\[
  \frac{1}{K}\cdot f(\phi_n^{-1}(C)) \leq i(C,f(\phi_n^{-1}(C)) \leq K\cdot f(\phi_n^{-1}(C)) 
\]
for all $n$ and all $f \in \mathcal{K}$. Taking minimum over $\lambda$, we get
\[
  m_n > \frac{1}{K} \cdot i(C,f(\phi_n^{-1}(C))
\]
for each $n$. Since $\phi_n$ are pairwise distinct, Lemma~\ref{lem:finitefilling}, we obtain $i(C,f(\phi_n^{-1}(C)) \to \infty$, and so does $m_n$, as $n \to \infty$.
\end{proof}


\section{Examples}

\subsection{Quasi-convex and convex}

\subsection{Geodesic currents as curve functionals}

\section{MCG action on functionals}

